\documentclass[10pt,aspectratio=169]{beamer}

\usetheme{Madrid}
\usecolortheme{dolphin}
\usefonttheme{professionalfonts}

\usepackage[T1]{fontenc}
\usepackage[utf8]{inputenc}
\usepackage{booktabs}
\usepackage{array}
\usepackage{hyperref}
\usepackage{xcolor}

\definecolor{codebg}{RGB}{245,245,245}
\definecolor{deepblue}{RGB}{0,70,130}
\setbeamercolor{title}{fg=deepblue}
\setbeamercolor{frametitle}{fg=deepblue}
\setbeamertemplate{navigation symbols}{}
\setbeamertemplate{itemize items}[circle]

\newcommand{\code}[1]{\texttt{\small\detokenize{#1}}}
\newcommand{\file}[1]{\texttt{\small\detokenize{#1}}}

\title[CHyLL v2 handoff]{Successful CHyLL v2 Examples: User Guide}
\subtitle{Data, models, and how to regenerate the downstream inputs}
\author{Kaito Iwasaki}
\date{Thursday walkthrough}

\begin{document}

\begin{frame}
  \titlepage
\end{frame}

\begin{frame}{Purpose of the Meeting}
  \begin{itemize}
    \item Give David and Bernardo a clean bundle with only successful examples.
    \item Map each file to the mathematical/computational object it contains.
    \item Show how to load the data immediately.
    \item Explain what the CHyLL v2 losses enforce and where they live in code.
    \item Show how to regenerate the exported lift data from the trained models.
    \item Identify what David needs for his downstream computations.
  \end{itemize}
\end{frame}

\begin{frame}{What We Are Sending}
  The file to send is:

  \vspace{0.5em}
  \begin{center}
    \file{2026-05-19_successful_examples_chyll_v2_repro.zip}
  \end{center}

  \vspace{0.5em}
  It contains:
  \begin{itemize}
    \item exported data for all successful examples,
    \item trained model weights and configs,
    \item minimal source code needed to reload/export,
    \item rollout figures for every successful model,
    \item training loss-history figures for every successful training,
    \item selected additional orientation figures,
    \item a README, examples table, requirements file, and hash manifest.
  \end{itemize}
\end{frame}

\begin{frame}{What Is Deliberately Not Included}
  The bundle excludes:
  \begin{itemize}
    \item smoke-test runs,
    \item failed compass uniform-initial-condition run,
    \item intermediate ablations,
    \item obsolete top-level \file{data/compass_gait/} files,
    \item older non-\file{chyll_v2} architecture outputs.
  \end{itemize}

  \vspace{0.5em}
  The goal is to avoid making David sort through bad or superseded data.
\end{frame}

\begin{frame}{The Successful Examples}
  \small
  \begin{tabular}{lll}
    \toprule
    Folder & System & Meaning \\
    \midrule
    \file{rimless_wheel_phaseB} & Rimless wheel & Stable limit cycle \\
    \file{bouncing_ball_phaseB_5impacts} & Bouncing ball & 5-impact lift \\
    \file{bouncing_ball_phaseB_15impacts} & Bouncing ball & 15-impact lift \\
    \file{compass_period2_4p75deg_phi1} & Compass gait & 4.75 deg, period 2 \\
    \file{compass_period4_5p00deg_phi2} & Compass gait & 5.00 deg, period 4 \\
    \file{compass_period8_5p02deg_phi3} & Compass gait & 5.02 deg, period 8 \\
    \file{compass_chaotic_cloud_5p20deg_phi4} & Compass gait & 5.20 deg, chaotic/cloud \\
    \bottomrule
  \end{tabular}
\end{frame}

\begin{frame}{One-Sentence Big Picture}
  Each example is a hybrid dynamical system trajectory embedded by CHyLL v2
  into a continuous latent space.

  \vspace{1em}
  The exported files give David:
  \begin{itemize}
    \item points along the latent trajectory,
    \item tangent directions along that trajectory,
    \item the trained model that generated them,
    \item and the existing cycling-signature output files.
  \end{itemize}
\end{frame}

\begin{frame}{Notation and Artifact Map}
  \small
  \begin{tabular}{ll}
    \toprule
    Artifact & Meaning \\
    \midrule
    \file{data/<case>/} & exported lifts and downstream outputs \\
    \file{models/<case>/} & trained CHyLL v2 checkpoint and run config \\
    \file{latent_lift_positions.npy} & sampled points in learned latent space \\
    \file{latent_lift_tangents.npy} & tangent vectors along the sampled lift \\
    \file{postimpact_*} & compass Poincare-section data after impact \\
    \file{cycling_signature_outputs/} & existing downstream outputs for comparison \\
    \bottomrule
  \end{tabular}

  \vspace{0.8em}
  Convention: \file{<case>} denotes one row of \file{examples.csv}.
\end{frame}

\begin{frame}[fragile]{Commands Used in the Demo}
  The only commands needed for the walkthrough are:

\begin{verbatim}
cd <bundle folder>
python -m pip install -r requirements.txt
python export_all_lifts.py
\end{verbatim}

  Replace \code{<bundle folder>} by the unzipped bundle path.
  If the path contains spaces, put quotes around it.
\end{frame}

\begin{frame}{Top-Level Folder Layout}
  After unzipping:

  \vspace{0.5em}
  \begin{tabular}{ll}
    \file{data/} & ready-to-use exported data \\
    \file{models/} & trained model weights and configs \\
    \file{source/} & source code needed to reload/export \\
    \file{figures/} & rollout, loss-history, and orientation plots \\
    \file{README.md} & text user guide \\
    \file{examples.csv} & table of examples \\
    \file{export_all_lifts.py} & regenerate exported lifts \\
    \file{requirements.txt} & Python package list \\
    \file{MANIFEST.csv} & file hashes \\
  \end{tabular}
\end{frame}

\begin{frame}{Rollout and Loss Figures}
  The \file{figures/} folder includes every successful rollout overlay:

  \begin{itemize}
    \item rimless wheel: short, index, and long rollout plots,
    \item bouncing ball: standard and index rollout plots,
    \item compass gait: one rollout plot for each incline case,
    \item loss-history plots for the six successful trainings,
    \item plus rank heatmaps and the compass latent-cluster panel.
  \end{itemize}

  \vspace{0.5em}
  These are diagnostic plots. The arrays for downstream computation are still
  the \file{data/<case>/latent_lift_*} files.
\end{frame}

\begin{frame}{Data Folder Anatomy}
  Each example folder in \file{data/} contains:

  \begin{itemize}
    \item \file{latent_lift_positions.npy}
    \item \file{latent_lift_tangents.npy}
    \item \file{latent_lift_positions.csv}
    \item \file{latent_lift_tangents.csv}
    \item \file{export_report.txt}
    \item \file{cycling_signature_outputs/}
  \end{itemize}

  \vspace{0.5em}
  The \file{.npy} files are best for Python. The \file{.csv} files are easier
  for spreadsheets and quick inspection.
\end{frame}

\begin{frame}{What Are Positions and Tangents?}
  \begin{itemize}
    \item \file{latent_lift_positions} is the sequence of points in latent space.
    \item \file{latent_lift_tangents} is the direction of motion at each point.
    \item The two arrays have the same shape.
    \item David's cycling-signature code uses both positions and tangents.
  \end{itemize}

  \vspace{0.5em}
  Example: a shape \code{(960, 7)} means:
  \begin{itemize}
    \item 960 sampled points,
    \item each point has 7 latent coordinates.
  \end{itemize}
\end{frame}

\begin{frame}{Shapes of the Exported Lifts}
  \small
  \begin{tabular}{ll}
    \toprule
    Example & Position/tangent shape \\
    \midrule
    Rimless wheel Phase-B & \code{(960, 7)} \\
    Bouncing ball, 5 impacts & \code{(1069, 7)} \\
    Bouncing ball, 15 impacts & \code{(2004, 7)} \\
    Compass 4.75 deg, period 2 & \code{(2032, 11)} \\
    Compass 5.00 deg, period 4 & \code{(2036, 11)} \\
    Compass 5.02 deg, period 8 & \code{(2036, 11)} \\
    Compass 5.20 deg, chaotic/cloud & \code{(2031, 11)} \\
    \bottomrule
  \end{tabular}
\end{frame}

\begin{frame}{Compass-Specific Files}
  The compass-gait folders also include:

  \begin{itemize}
    \item \file{postimpact_states.npy}: physical post-impact states.
    \item \file{postimpact_latent_z.npy}: encoded post-impact section points.
    \item \file{postimpact_labels.npy}: cluster labels for the period diagnostic.
  \end{itemize}

  \vspace{0.5em}
  These files support the incline-cascade interpretation:
  period 2, period 4, period 8, and chaotic/cloud.
\end{frame}

\begin{frame}{Model Folder Anatomy}
  Each folder in \file{models/} contains:

  \begin{itemize}
    \item \file{model.pt}: learned neural-network weights.
    \item \file{config.json}: model/system settings used for that run.
    \item \file{train_log.jsonl}: training log.
    \item optional \file{console.*.log}: raw console output for some runs.
  \end{itemize}

  \vspace{0.5em}
  The data can be used without loading the models. The models are included so
  David can regenerate the data or modify export settings.
\end{frame}

\begin{frame}{Learned Objects}
  CHyLL v2 learns three maps on the augmented mapping-cylinder state
  \(y=(x,s)\in X'\):

  \vspace{0.5em}
  \begin{align*}
    E_\theta &: X' \to \mathbb{R}^m && \text{encoder},\\
    V_\theta &: \mathbb{R}^m \to \mathbb{R}^m && \text{latent vector field},\\
    D_\theta &: \mathbb{R}^m \to X' && \text{decoder}.
  \end{align*}

  \vspace{0.3em}
  Here \(m=7\) for rimless/bouncing and \(m=11\) for compass gait.
\end{frame}

\begin{frame}{Implementation Map}
  \small
  \begin{tabular}{ll}
    \toprule
    Mathematical object & Code location \\
    \midrule
    \(E_\theta, V_\theta, D_\theta\) & \file{chyll_v2/networks.py} \\
    latent ODE rollout & \file{chyll_v2/ode.py} \\
    loss terms \(L_x,L_z,L_g,L_v,L_c\) & \file{chyll_v2/losses.py} \\
    training loop / checkpoint saving & \file{chyll_v2/train.py} \\
    system simulators and reset maps & \file{chyll_v2/systems/} \\
    export scripts & \file{cycling_signature/export/} \\
    \bottomrule
  \end{tabular}

  \vspace{0.5em}
  All paths here are relative to \file{source/chyll_v2/}.
\end{frame}

\begin{frame}{Forward Pass During Training}
  For a sampled trajectory slice \(y_0,\ldots,y_{T-1}\):

  \vspace{0.5em}
  \begin{align*}
    z_i &= E_\theta(y_i),\\
    \widehat z_i &= \Psi_\theta(t_i; z_0), \qquad \dot z=V_\theta(z),\\
    \widehat y_i &= D_\theta(\widehat z_i).
  \end{align*}

  \vspace{0.3em}
  In code this is \code{forward_step(...)} in \file{train.py}: encode the
  observed slice, integrate the Neural ODE, decode the integrated latent
  trajectory, then call \code{compute_losses(...)}.
\end{frame}

\begin{frame}{Composite Loss}
  The optimized objective is:

  \vspace{0.5em}
  \[
    L = w_x L_x + w_z L_z + w_g L_g + w_v L_v + w_c L_c.
  \]

  \vspace{0.5em}
  Included successful examples use:
  \begin{itemize}
    \item rimless/bouncing: \((w_x,w_z,w_g,w_v,w_c)=(1,1,3,1,1)\),
    \item compass: \((w_x,w_z,w_g,w_v,w_c)=(1,1,3,0,1)\).
  \end{itemize}

  The exact values are in each \file{models/<case>/config.json}.
\end{frame}

\begin{frame}{What the Losses Enforce}
  \scriptsize
  \begin{tabular}{lll}
    \toprule
    Term & Mathematical role & Code helper \\
    \midrule
    \(L_x\) & \(D_\theta(\widehat z_i)\approx y_i\) reconstruction & \code{reconstruction_loss} \\
    \(L_z\) & \(\Psi_\theta(t_i,E_\theta(y_0))\approx E_\theta(y_i)\) & \code{latent_consistency_loss} \\
    \(L_g\) & contract reset/gluing pairs in latent space & \code{gluing_loss} \\
    \(L_v\) & match pre/post seam velocities where enabled & \code{seam_velocity_loss} \\
    \(L_c\) & keep latent coordinates from collapsing & \code{anti_collapse_loss} \\
    \bottomrule
  \end{tabular}

  \vspace{0.7em}
  \normalsize
  The glue/seam masks are data-driven: they are inferred from the sampled
  mapping-cylinder coordinate \(s\), not from symbolic guard equations.
\end{frame}

\begin{frame}{How the Learned Models Are Found}
  \begin{itemize}
    \item Training writes one run directory per successful case.
    \item \file{model.pt} is the PyTorch \code{state_dict}.
    \item \file{config.json} reconstructs dimensions, system parameters, loss weights,
          and ODE/training settings.
    \item \file{train_log.jsonl} records \(L_x,L_z,L_g,L_v,L_c\) over training.
  \end{itemize}

  \vspace{0.5em}
  The handoff copies only the successful checkpoints into \file{models/}.
\end{frame}

\begin{frame}[fragile]{How the Export Script Uses a Model}
  Each entry in \file{export_all_lifts.py} specifies:

\begin{verbatim}
"model": ROOT / "models" / <case> / "model.pt"
"config": ROOT / "models" / <case> / "config.json"
"script": ... / "prepare_*_lift.py"
\end{verbatim}

  The export script rebuilds the network from \file{config.json}, loads
  \file{model.pt}, simulates the corresponding hybrid example, encodes it
  with \(E_\theta\), computes tangents, and writes \file{regenerated/<case>/}.
\end{frame}

\begin{frame}{Three Levels of Use}
  \begin{enumerate}
    \item \textbf{Use existing data.}
      Load \file{latent_lift_positions.npy} and \file{latent_lift_tangents.npy}.
    \item \textbf{Regenerate lift data.}
      Run \file{python export_all_lifts.py}.
    \item \textbf{Inspect or modify model-derived exports.}
      Edit only export parameters first; retraining is not needed for Thursday.
  \end{enumerate}

  \vspace{0.5em}
  For Thursday, we should demo levels 1 and 2 and show where the model files
  enter the regeneration step.
\end{frame}

\begin{frame}[fragile]{Loading Data in Python}
  Minimal example:

\begin{verbatim}
import numpy as np

case = "compass_period4_5p00deg_phi2"
Z = np.load(f"data/{case}/latent_lift_positions.npy")
T = np.load(f"data/{case}/latent_lift_tangents.npy")

print(Z.shape)
print(T.shape)
\end{verbatim}

  Here \code{Z} is the latent position array and \code{T} is the tangent array.
\end{frame}

\begin{frame}[fragile]{Loading Compass Post-Impact Data}
  For compass-gait period/cascade analysis:

\begin{verbatim}
states = np.load(f"data/{case}/postimpact_states.npy")
section_z = np.load(f"data/{case}/postimpact_latent_z.npy")
labels = np.load(f"data/{case}/postimpact_labels.npy")

print(states.shape)
print(section_z.shape)
print(labels.shape)
\end{verbatim}

  The labels are the cluster IDs used to count period-2, period-4, period-8,
  or many/cloud behavior.
\end{frame}

\begin{frame}[fragile]{Loading CSV Files Instead}
  If someone prefers CSV:

\begin{verbatim}
import numpy as np

Z = np.loadtxt("data/rimless_wheel_phaseB/latent_lift_positions.csv")
T = np.loadtxt("data/rimless_wheel_phaseB/latent_lift_tangents.csv")

print(Z.shape)
\end{verbatim}

  The CSV files are space-delimited; use a space delimiter in any non-Python
  tool.
\end{frame}

\begin{frame}{What the Export Script Does}
  The script:
  \begin{itemize}
    \item loads each bundled \file{model.pt},
    \item reads the matching \file{config.json},
    \item simulates the appropriate hybrid trajectory,
    \item encodes it into CHyLL v2 latent space,
    \item computes tangent directions,
    \item writes outputs to \file{regenerated/<case>/}.
  \end{itemize}

  \vspace{0.5em}
  It does not retrain models. It only regenerates exported data.
\end{frame}

\begin{frame}[fragile]{Python Setup}
  From the unzipped bundle folder:

\begin{verbatim}
python -m pip install -r requirements.txt
\end{verbatim}

  Then:

\begin{verbatim}
python export_all_lifts.py
\end{verbatim}

  This should create:

\begin{verbatim}
regenerated/rimless_wheel_phaseB/
regenerated/bouncing_ball_phaseB_5impacts/
...
\end{verbatim}
\end{frame}

\begin{frame}{The Role of \file{source/}}
  \file{source/} contains the minimal CHyLL v2 code needed to reload models:

  \begin{itemize}
    \item network architecture,
    \item system simulators,
    \item lift exporters,
    \item downstream helper scripts.
  \end{itemize}

  \vspace{0.5em}
  It is included so David does not need to reconstruct the repo structure from
  scratch just to regenerate the data.
\end{frame}

\begin{frame}{Rimless and Bouncing-Ball Tangents}
  For rimless wheel and bouncing ball:

  \begin{itemize}
    \item exported tangents use tag-aware finite differences,
    \item this differentiates within each arc/bridge piece,
    \item it avoids artificial chords across concatenation boundaries.
  \end{itemize}

  \vspace{0.5em}
  In the export script this corresponds to:

  \begin{center}
    \code{--tangent-source diff --tangent-mode tagaware}
  \end{center}
\end{frame}

\begin{frame}{Compass Tangents}
  For the compass incline cascade:

  \begin{itemize}
    \item exported tangents use the learned latent vector field,
    \item this matches the included files,
    \item it is the tangent convention used for the current compass cascade data.
  \end{itemize}

  \vspace{0.5em}
  In the export script this corresponds to:

  \begin{center}
    \code{--tangent-source vfield}
  \end{center}
\end{frame}

\begin{frame}{Downstream Computation Boundary}
  Existing outputs are already included under:

  \vspace{0.5em}
  \begin{center}
    \file{data/<case>/cycling_signature_outputs/}
  \end{center}

  \vspace{0.5em}
  \begin{itemize}
    \item The handoff responsibility is to provide clean positions, tangents,
          and labels.
    \item David can plug those arrays into his existing downstream workflow.
    \item We do not need to explain or debug his downstream environment in the
          main walkthrough.
  \end{itemize}
\end{frame}

\begin{frame}{Common Mistakes}
  \begin{itemize}
    \item Running commands from the wrong folder.
    \item Forgetting quotes around paths with spaces.
    \item Mixing old \file{data/compass_gait} files with CHyLL v2 files.
    \item Using the wrong model with the wrong config.
    \item Expecting \file{export_all_lifts.py} to train models.
  \end{itemize}
\end{frame}

\begin{frame}{Sanity Checks}
  After loading data:

  \begin{itemize}
    \item positions and tangents should have the same shape,
    \item tangent norms should be close to 1,
    \item compass post-impact states should be \code{(256, 4)},
    \item compass post-impact latent points should be \code{(256, 11)}.
  \end{itemize}

  \vspace{0.5em}
  These checks are already done for the shipped bundle.
\end{frame}

\begin{frame}[fragile]{Tiny Python Sanity Check}
\begin{verbatim}
import numpy as np

case = "rimless_wheel_phaseB"
Z = np.load(f"data/{case}/latent_lift_positions.npy")
T = np.load(f"data/{case}/latent_lift_tangents.npy")

norms = np.linalg.norm(T, axis=1)
print(Z.shape, T.shape)
print(norms.min(), norms.max())
\end{verbatim}

  The min and max should both print very close to 1.
\end{frame}

\begin{frame}{Suggested Thursday Agenda}
  \begin{enumerate}
    \item Confirm everyone is using the successful-examples zip.
    \item Walk through \file{examples.csv} and folder layout.
    \item Load one \file{.npy} data file in Python.
    \item Explain the learned maps and loss terms at the level of the code.
    \item Regenerate one example using \file{export_all_lifts.py}.
    \item Explain the cycling-signature output folders.
  \end{enumerate}
\end{frame}

\begin{frame}{Live Demo Plan}
  Recommended demo order:

  \begin{enumerate}
    \item Open the zip folder.
    \item Open \file{README.md}.
    \item Open \file{examples.csv}.
    \item Load \file{data/compass_period4_5p00deg_phi2/latent_lift_positions.npy}.
    \item Run \file{python export_all_lifts.py}.
    \item Compare \file{data/} and \file{regenerated/}.
  \end{enumerate}
\end{frame}

\begin{frame}{What David Can Safely Modify}
  \begin{itemize}
    \item Which example/case to process.
    \item Cycling-signature parameters.
    \item Number and length of random subsegments.
    \item Export destination folder.
    \item Whether to use existing data or regenerated data.
  \end{itemize}

  \vspace{0.5em}
  He should not need to retrain models for the immediate downstream task.
\end{frame}

\begin{frame}{Best Follow-Up Deliverable}
  After the meeting, it would be useful to send:

  \begin{itemize}
    \item the final zip,
    \item this slide deck,
    \item a one-paragraph note saying which examples are successful,
    \item and the exact command David should run first.
  \end{itemize}

  \vspace{0.5em}
  First command:
  \begin{center}
    \code{python export_all_lifts.py}
  \end{center}
\end{frame}

\begin{frame}{Summary}
  \begin{itemize}
    \item The zip is a clean successful-examples bundle.
    \item Data can be used immediately from \file{data/}.
    \item Models and source are included for regeneration.
    \item \file{export_all_lifts.py} is the easiest entry point.
    \item Downstream outputs are included for comparison; David can rerun them
          in his own workflow.
  \end{itemize}
\end{frame}

\end{document}
